Una de las razones principales del éxito en producción de estas
plataformas no reside únicamente en la calidad de sus interfaces de
corrección, sino en su acoplamiento con la infraestructura informática del
centro. La sincronización automática de listas desde el \ac{lms}, la
entrega de calificaciones al \textit{gradebook}, el inicio de sesión
institucional mediante \acs{sso} y la conformidad con estándares como
\acs{lti} 1.3 reducen drásticamente la fricción de adopción. En el
contexto español, donde Moodle es el \acs{lms} predominante en
universidades públicas, las plataformas que integran de forma nativa con
Moodle parten con una clara ventaja operativa.
